\begin{proposition}[A convergent arithmetic refinement of the far inverse]
\label{prop:v120-inverse-polynomial}
If $D_g\preceq U\preceq mD_g$, $D_g\succeq\gamma I>0$,
$m>1$, let $A=D_g^{-1/2}UD_g^{-1/2}$ denote the bounded
operator represented by the preconditioned form, and set
\[
 S=I-A/m,\qquad
 P_k(A)=\frac1m\sum_{j=0}^{k-1}S^j+S^k,\quad k\ge0.
\]
The empty sum gives $P_0=I$. Then
\begin{equation}\label{eq:v120-inverse-polynomial}
 A^{-1}\preceq P_{k+1}(A)\preceq P_k(A)\preceq I,
 \qquad
 0\preceq P_k(A)-A^{-1}\preceq(1-1/m)^k I.
\end{equation}
Thus $D_g^{-1/2}P_k(A)D_g^{-1/2}$ are decreasing upper
bounds for $U^{-1}$, with error at most
$(1-1/m)^kD_g^{-1}$ in form order. In particular,
\begin{equation}\label{eq:v120-first-inverse-refinement}
 U^{-1}\preceq D_g^{-1}
       -m^{-1}D_g^{-1}(U-D_g)D_g^{-1}.
\end{equation}
Products involving $U$ in this formula are understood through the
bounded preconditioned form.

For the actual far operators used in
Proposition~\ref{prop:v120-lambda4-tail}, including their shift
$\mathsf W_4-10^{-8}D$, one can take
\[
 m=67\quad\hbox{on the even tail }n>512,\qquad
 m=335\quad\hbox{on the odd tail }n>1536.
\]
These bounds cover all omitted Fourier modes at this fixed window.
\end{proposition}
\begin{proof}
The spectral theorem gives $I\preceq A\preceq mI$ and
$0\preceq S\preceq(1-1/m)I$. The geometric identity is
\[
 P_k(A)-A^{-1}=S^k(I-A^{-1}).
\]
Its factors commute and are positive. This proves the error bound;
$P_k-P_{k+1}=m^{-1}S^k(A-I)\succeq0$ proves monotonicity.
Conjugate by $D_g^{-1/2}$, using the inverse identity for the
positive closed form $U$. The case $k=1$ is
\eqref{eq:v120-first-inverse-refinement}.

To verify the two concrete constants, first strengthen
\eqref{eq:v115-diagonal-sharp} to
\begin{equation}\label{eq:v120-two-sided-arch}
 |a_n-\log(|n|/L)|\le E_L(|t_n|),\qquad n\ne0.
\end{equation}
For the upper bound use the same Binet formula at
$w=5/4+it/2$. It gives
\[
 \Re\log w-\log(t/2)
 =\tfrac12\log(1+25/(4t^2))\le25/(8t^2),
\]
while $-\Re(1/(2w))-\Re(1/z)\le0$, $z=1/4+it/2$.
The Binet remainder is at most $1/(15t)$ in absolute value.
Because $25/8<7/2$, the existing absolute trigamma and exponential
series bounds give the same $E_L$ for the upper error.

Write $\mathsf W_4=D+R_0$. From
\eqref{eq:v114-b-d}--\eqref{eq:v114-b-bound}, a valid bound is
\[
 \norm{R_0}\le C_R:=2\pi B_*+32h_L/L+2P_4.
\]
Here $P_4$ may omit the zero full-length translation $m=16$.
The pole diagonal bound is $32h_L/L$; its factor $L$ is essential.
The commutator bound is taken before the parity compression, so it
also bounds either normalized parity sector. With $c=10^{-8}$,
$g_n=(1-c)\{\log(n/L)-E_L(t_n)\}-\kappa$, and
$t_*=2\pi(N+1)/L$, the already proved lower bound and
\eqref{eq:v120-two-sided-arch} imply
\[
 D_g\preceq U\preceq
 D_g+\{\kappa+2(1-c)E_L(t_*)+C_R\}I.
\]
Since $g_n$ is positive and increasing, it suffices that
\begin{equation}\label{eq:v120-inverse-m-bound}
 m\ge1+\frac{\kappa+2(1-c)E_L(t_*)+C_R}{g_{N+1}}.
\end{equation}
Outward interval evaluation gives $C_R<34.450462$. The right side
of \eqref{eq:v120-inverse-m-bound} is below $66.76640$ for the
even $N=512$ choice, and below $334.71708$ for the odd $N=1536$
choice. The same prime bound, diagonal errors and inverse weights
as the tail certificate are used. Strict interval checks at two
precisions verify the stated integers.
\end{proof}

For a matrix of residuals $Y$, the error of the $k$th inverse bound
is at most
\[
 (1-1/m)^k Y^*D_g^{-1}Y.
\]
Along growing windows this part of the error is therefore controlled
if $\norm{D_{g,\lambda}^{-1/2}Y_\lambda}^2
 e^{-k_\lambda/m_\lambda}\to0$, with the actual values of
$m_\lambda$ retained. The explicit integers above apply only at
$\lambda=4$. Evaluating the signed powers and their complete remote
contributions, and proving the remaining low-head sign, are separate
obligations. The inverse iteration proves neither of those signs.

